A mass is subject to a damping force proportional to its velocity, which means that the equation of motion takes the form $\ddot{x} = -A\dot{x}$, where $A$ is some constant. If the initial speed is 2 m/s, and if $A = 1 \text{ s}^{-1}$, how far has the mass traveled at 1 s? 10 s? 100 s? You should find that the distance approaches a limiting value.

Now assume that the mass is subject to a damping force proportional to the square of its velocity, which means that the equation of motion now takes the form $\ddot{x} = -A\dot{x}^2$, where $A$ is some constant. If the initial speed is 2 m/s, and if $A = 1 \text{ m}^{-1}$, how far has the mass traveled at 1 s? 10 s? 100 s? How about some larger powers of 10? You should find that the distance keeps growing, but slowly like the log of $t$. (The results for these two forms of the damping are consistent with the results of Problem 1.5.)